%! Characteristics of Logarithmic Functions — Unit 8, Lesson 8.0 — Homework
\documentclass[10pt]{article}
\usepackage{algebra2-article}
\usepackage{algebra2-boxes}
\newcommand{\TallMath}[1]{$\displaystyle #1\rule[-1.4em]{0pt}{3.2em}$}
\newcommand{\lgrid}[4]{%
  \draw[step=1, linegray!40, very thin] (#1,#3) grid (#2,#4);
  \draw[->, semithick, charcoal] (#1,0)--(#2,0) node[below right, font=\scriptsize]{$x$};
  \draw[->, semithick, charcoal] (0,#3)--(0,#4) node[left, font=\scriptsize]{$y$};}
% pgfmath has no log_b, so every logarithmic curve is drawn by parameterizing on y:
%   y = log_b(x-h) + k   <=>   x = b^(y-k) + h.   Args: {b}{h}{k}{ymin}{ymax}
\newcommand{\logcurve}[5]{\draw[very thick, forest, domain=#4:#5, samples=120, smooth]
  plot ({pow(#1,\x-(#3))+(#2)},{\x});}

\begin{document}

\pageheader{Unit 8, Lesson 8.0}{Homework}
\namedateperiod

\vspace{0.10in}

\begin{notesbox}{Practice}
Show how you read each feature. Write every domain and range in \textbf{interval notation}, and
remember that a logarithm's boundary is never included.

\begin{enumerate}[leftmargin=1.9em, itemsep=12pt]
  \item \textbf{From the equation only.} Use the argument --- no graphing.
    \begin{center}
    \renewcommand{\arraystretch}{1.6}
    \begin{tabularx}{\linewidth}{@{}l X X X X@{}}
    \hline
    \textbf{Function} & \textbf{Condition on the argument} & \textbf{Domain} & \textbf{Vertical asymptote} & \textbf{Runs left or right?} \\
    \hline
    $\log_3(x-2)$ & \blank{1.8cm} & \blank{2.0cm} & \blank{1.5cm} & \blank{1.5cm} \\
    $\log(x+7)$ & \blank{1.8cm} & \blank{2.0cm} & \blank{1.5cm} & \blank{1.5cm} \\
    $\log_2(4x)$ & \blank{1.8cm} & \blank{2.0cm} & \blank{1.5cm} & \blank{1.5cm} \\
    $\log_5(10-2x)$ & \blank{1.8cm} & \blank{2.0cm} & \blank{1.5cm} & \blank{1.5cm} \\
    \hline
    \end{tabularx}
    \end{center}

  \item \textbf{A shift that moves the intercept but not the asymptote.} The graph is
        $m(x) = \log_2 x - 3$.
    \begin{center}
    \begin{tikzpicture}[scale=0.44]
      \lgrid{-2}{13}{-6}{2}
      \draw[navy, dashed, semithick] (0,-6)--(0,2);
      \begin{scope}
        \clip (-2,-6) rectangle (13,2);
        \logcurve{2}{0}{-3}{-6}{0.70}
      \end{scope}
      \fill[charcoal] (1,-3) circle (3pt) node[below right, font=\scriptsize]{$(1,-3)$};
      \fill[charcoal] (2,-2) circle (3pt);
      \fill[charcoal] (4,-1) circle (3pt) node[above left, font=\scriptsize]{$(4,-1)$};
      \fill[charcoal] (8,0) circle (3pt) node[above left, font=\scriptsize]{$(8,0)$};
    \end{tikzpicture}
    \end{center}
    Domain: \blank{3.0cm} \quad Range: \blank{3.0cm} \quad Vertical asymptote: \blank{2.0cm} \\[3pt]
    $x$-intercept: \blank{1.8cm} \quad $y$-intercept: \blank{1.8cm} \\[3pt]
    Increasing or decreasing on its domain? \blank{2.4cm} \quad Absolute extrema: \blank{2.0cm} \\[3pt]
    $m(x)<0$ on \blank{2.6cm}; \quad $m(x)>0$ on \blank{2.6cm} \\[3pt]
    End behavior: as $x\to+\infty$, $m(x)\to\blank{1.4cm}$. \; As $x\to\blank{1.4cm}$,
    $m(x)\to\blank{1.4cm}$ \\[3pt]
    The parent $y=\log_2 x$ has $x$-intercept $(1,0)$ and asymptote $x=0$. Subtracting $3$ moved
    \emph{one} of those two. Which one, and why not the other? \blank{5.0cm}

  \item \textbf{A base between 0 and 1.} The graph is $n(x) = \log_{1/3} x$.
    \begin{center}
    \begin{tikzpicture}[scale=0.44]
      \lgrid{-2}{11}{-3}{3}
      \draw[navy, dashed, semithick] (0,-3)--(0,3);
      \begin{scope}
        \clip (-2,-3) rectangle (11,3);
        \logcurve{0.333333}{0}{0}{-2.19}{3}
      \end{scope}
      \fill[charcoal] (1,0) circle (3pt) node[above right, font=\scriptsize]{$(1,0)$};
      \fill[charcoal] (3,-1) circle (3pt) node[below right, font=\scriptsize]{$(3,-1)$};
      \fill[charcoal] (9,-2) circle (3pt) node[below right, font=\scriptsize]{$(9,-2)$};
    \end{tikzpicture}
    \end{center}
    Domain: \blank{3.0cm} \quad Range: \blank{3.0cm} \quad Vertical asymptote: \blank{2.0cm} \\[3pt]
    $x$-intercept: \blank{1.8cm} \quad $y$-intercept: \blank{1.8cm} \\[3pt]
    Increasing or decreasing? \blank{2.4cm} \quad Absolute extrema: \blank{2.4cm} \\[3pt]
    End behavior: as $x\to+\infty$, $n(x)\to\blank{1.2cm}$; \; as $x\to\blank{1.4cm}$,
    $n(x)\to\blank{1.2cm}$ \\[3pt]
    Compared with $y=\log_3 x$ from class, list every feature that is the \textbf{same}:
    \blank{5.0cm}

  \item \textbf{Explain.} Every parent logarithmic graph $y=\log_b x$ passes through $(1,0)$ and
        never crosses the $y$-axis --- no matter what the base is. Explain both facts, using the
        words \textbf{exponent} and \textbf{asymptote}.
        \writelines{3}
\end{enumerate}
\end{notesbox}

\begin{extensionbox}
\textbf{Building a log graph from an exponential table.} Below is a table for $y=3^x$. You will not
need a calculator.
\begin{center}
\renewcommand{\arraystretch}{1.25}
\begin{tabular}{|c|c|c|c|c|c|}
\hline
$x$ & $-2$ & $-1$ & $0$ & $1$ & $2$ \\ \hline
$y=3^x$ & $\tfrac19$ & $\tfrac13$ & $1$ & $3$ & $9$ \\ \hline
\end{tabular}
\end{center}
\begin{enumerate}[label=(\alph*), leftmargin=1.8em, itemsep=5pt]
  \item Trade the rows to build a table for $y=\log_3 x$. Which row becomes the inputs?
  \item Using only your new table, state the domain and range of $y=\log_3 x$. Explain how each one
        came from a feature of $y=3^x$.
  \item The graph of $y=3^x$ has a horizontal asymptote at $y=0$. Reflect that line over $y=x$: what
        line do you get, and what is it called on the logarithmic graph?
  \item $y=3^x$ has a $y$-intercept at $(0,1)$ and no $x$-intercept at all. Use the coordinate swap
        to predict both intercept facts for $y=\log_3 x$ --- then check them against your table.
  \item In Unit~6, $y=x^2$ had to be cut down to $x\ge 0$ before $\sqrt{x}$ could undo it, but
        $y=3^x$ needs no such cut. What is true about an exponential graph that makes the difference?
\end{enumerate}
\writelines{5}
\end{extensionbox}

\begin{spiralbox}
\textbf{Coming next --- Lesson 8.1: Introduction to Logarithms.} Today $\log_b x$ was the
\emph{name of a curve}: you read it off a graph and used it to decide which inputs were legal. Next
you will learn to \emph{compute} with it --- writing $\log_b x = y$ and $b^y = x$ as one statement
read two ways, evaluating logs by asking ``$b$ to what power?'', meeting the common log, and using
change of base so a calculator can handle any base at all. That is also where today's domain rule
gets its full algebraic argument.
\end{spiralbox}

\end{document}
